\documentclass[10pt,a4paper]{amsart}
\usepackage[margin=19mm]{geometry}
\usepackage{amsmath,amssymb,mathtools}
\usepackage[scr=rsfs,cal=euler]{mathalfa}
\usepackage{xfrac}
\usepackage[unicode,hidelinks,bookmarks=false]{hyperref}
\newcommand{\ie}{i.e.\ }
\newcommand{\cf}{cf.\ }
\newcommand{\resp}{respectively\ }
\newcommand{\Subsection}{\subsection}
\providecommand{\nobreakdash}{\nobreak\hbox{-}}
\setlength{\emergencystretch}{2em}
\multlinegap0pt
\usepackage{bm}
\usepackage{xcolor}
\usepackage{mathtools}
\usepackage{amssymb}
\usepackage{tikz}
\usepackage{esint}
\usepackage{enumerate}
\usepackage{caption}
\newtheorem{thm}{Theorem}
\newcommand{\Cth}{\cC_\theta}
\newcommand{\Si}{\Sigma}
\newcommand{\bbB}{\mathbb{B}}
\newcommand{\bbR}{\mathbb{R}}
\newcommand{\bbS}{\mathbb{S}}
\newcommand{\cC}{\mathcal{C}}
\newcommand{\eq}[1]{\begin{equation}\begin{alignedat}{2}#1\end{alignedat}\end{equation}}
  \def\in{in }%
\newcommand{\ip}[2]{\langle #1,#2 \rangle}
\newcommand{\op}[1]{\operatorname{#1}}
\newcommand{\ov}{\overline}
  \def\theta{theta}%
\newcommand{\wh}{\widehat}
\title{Capillary $L_p$-Christoffel-Minkowski problem}
\author{Yingxiang Hu, Mohammad N. Ivaki}
\date{Source: arXiv:2512.15464v2 (CC BY 4.0)}
\begin{document}
\maketitle

\noindent\emph{Source and licence.} Capillary $L_p$-Christoffel-Minkowski problem. Version arXiv:2512.15464v2 (CC BY 4.0), distributed by arXiv under the Creative Commons Attribution 4.0 International licence, http://creativecommons.org/licenses/by/4.0/, as recorded in the arXiv licence metadata for this exact version and shown on the abstract page https://arxiv.org/abs/2512.15464v2. Full licence notice: \texttt{licenses/arxiv-2512.15464-CC-BY-4.0.txt}.\par\smallskip

\noindent\textit{Editorial scope.} Counted unit: the article's thm body (source0.tex lines 437--448). The article's complete original proof of it is delivered with it (source0.tex lines 450--619, one \texttt{proof} environment of 170 lines). No mathematics is written, repaired or re-derived: the statement and the proof are the article's own lines, in the article's own notation, and no retained line is substituted or re-flowed.  No further source-local context is needed: every cross-reference of every retained line resolves inside the two delivered spans.  Nothing was omitted from any retained span, so the package carries no omission record.  No retained line contains a citation, so the wrapper carries no bibliography and no bibliography entry.  No retained line contains a figure environment, an image-inclusion command or drawing code, so no image is delivered and the package declares no asset dependency.  Editorial formatting only, in addition: an AMS \texttt{amsart} preamble that restates the article's own declarations and macros used by the retained lines, namely the environments \texttt{thm} on separate counters, together with the standard packages those lines need.  The retained environments print in this document as Theorem 1 in order of appearance, whereas the article prints the same items with the numbering of its own full document.  The complete native source of the article as distributed in the arXiv e-print for this exact version is bundled unchanged as \texttt{sources/191177/source0.tex}.
\begin{thm}\label{lem:tilde-nu-parallel body}
Let $\Sigma\subset\ov{\bbR^{n+1}_+}$ be a strictly convex, $\theta$-capillary hypersurface.
For $t>0$ define
\eq{ \phi_t:\Sigma\to\ov{\bbR^{n+1}_+},\quad \phi_t(x):=x+t \tilde\nu(x).
}
Then $\Sigma_t:=\phi_t(\Sigma)$ is a strictly convex, $\theta$-capillary hypersurface. Moreover, the (standard) outer parallel convex body
\eq{
  K_t:=\bigl(\widehat{\Sigma}-t\cos\theta\,e_{n+1}\bigr)+t\,\bbB,
}
and the capillary outer parallel convex body are related via $\wh{\Si_t}=K_t\cap \ov{\bbR^{n+1}_+}$.
In addition, we have $\Si_t=\Si+t\Cth$.
\end{thm}

\begin{proof}
Let $P=\{x_{n+1}=0\}$. Define $f(x)=\langle x,e_{n+1}\rangle=x_{n+1}$ and
\eq{
  g(x)=\langle \tilde\nu(x),e_{n+1}\rangle
       =\langle \nu(x),e_{n+1}\rangle-\cos\theta.
}
Then
\eq{\label{eq:height-after-shift}
  \langle \phi_t(x),e_{n+1}\rangle = f(x)+t\,g(x).
}   

\emph{Step 1.}
On $\partial\Sigma$ we have $f=0$ (since $\partial\Sigma\subset P$) and $g=0$ (by capillarity),
hence $(f+tg)(x)=0$ for $x\in\partial\Sigma$, i.e. $\phi_t(\partial\Sigma)\subset P$. Moreover, since $\nu(\op{int}(\Si))\subset \op{int}(\bbS^n_{\theta})$, $f+t g>0$ on $\op{int}(\Sigma)$ for any $t>0$. 

\emph{Step 2.}
Let $x\in\Sigma$ and choose an orthonormal basis $\{e_1,\dots,e_n\}$ of $T_x\Sigma$
consisting of principal directions, so that
\eq{
  {}^{\Si}\nabla_{e_i}\nu=\kappa_i e_i,\quad i=1,\dots,n,
}
with principal curvatures $\kappa_i$.
We have ${}^{\Si}\nabla_{e_i}\tilde\nu={}^{\Si}\nabla_{e_i}\nu$, and therefore
\eq{
  d\phi_t(e_i)=e_i+t\,{}^{\Si}\nabla_{e_i}\tilde\nu
  =e_i+t\,{}^{\Si}\nabla_{e_i}\nu
  =(1+t\kappa_i)e_i.
}
Thus $d\phi_t(T_x\Sigma)$ is spanned by $\{e_1,\dots,e_n\}$, $\phi_t$ is a smooth immersion, and the oriented unit normal of $\Sigma_t=\phi_t(\Sigma)$
at $y=\phi_t(x)$ equals $\nu(x)$, i.e.
\eq{
  \nu_t(y)=\nu(x)\quad\text{for }y=\phi_t(x).
}
Next we show that $\phi_t$ is an injective immersion and thus an embedding. Assume $\phi_t(x)=\phi_t(x')$ for some
$x,x'\in\Sigma$. Then
\eq{
  x+t(\nu(x)-\cos\theta\,e_{n+1})=x'+t(\nu(x')-\cos\theta\,e_{n+1}),
}
hence
\eq{\label{eq:inj-start}
  x-x'=t\bigl(\nu(x')-\nu(x)\bigr).
}
Taking the inner product with $\nu(x)$ gives
\eq{\label{eq:inj-ip}
  \ip{x-x'}{\nu(x)}=t\bigl(\ip{\nu(x')}{\nu(x)}-1\bigr).
}
Since $\wh{\Si}$ is convex and $\nu(x)$ is the outer normal vector to $\Si$ at $x$,
\eq{
  \ip{x'-x}{\nu(x)}\le 0.
}
On the other hand, $\ip{\nu(x')}{\nu(x)}\le 1$, hence by \eqref{eq:inj-ip}
\eq{
  \ip{x-x'}{\nu(x)}=0
  \quad\text{and}\quad
  \ip{\nu(x')}{\nu(x)}=1.
}
Thus $\nu(x')=\nu(x)$. Since $\Sigma$ is strictly convex, the Gauss map
$\nu:\Sigma\to\bbS_{\theta}^n$ is injective, hence $x'=x$.

\emph{Step 3.}
If $y=\phi_t(x)$ with $x\in\partial\Sigma$, then by step 1 and step 2 we have $y\in P$ and
$\nu_t(y)=\nu(x)$. Therefore,
\eq{
  \langle \nu_t(y),e_{n+1}\rangle
  = \langle \nu(x),e_{n+1}\rangle
  = \cos\theta,
}
so $\Sigma_t$ meets $P$ with the same contact angle $\theta$.


\emph{Step 4.}
Let $x\in\Sigma$ and set
\eq{
y:=\phi_t(x)=x+t(\nu(x)-\cos\theta\,e_{n+1}).
}
We claim that $y\in\partial K_t$ and that $\nu(x)$ is an outer normal of $K_t$ at $y$.

Note that $y\in K_t$. Since $\nu(x)$ is an outer unit normal of the convex body $\widehat{\Sigma}$ at $x$, we have
\eq{\label{eq:support-K}
  \ip{z-x}{\nu(x)}\le 0\quad\forall\,z\in \widehat{\Sigma}.
}
Let $w\in K_t$. Then for some $z\in \widehat{\Sigma}$ and $b\in\bbB$:
\eq{
  w=(z-t\cos\theta\,e_{n+1})+tb.
}
Moreover, we have
\eq{
  \ip{w-y}{\nu(x)} = \ip{z-x}{\nu(x)} + t\ip{b}{\nu(x)} - t.
}
Using \eqref{eq:support-K}, we obtain
\eq{
  \ip{w-y}{\nu(x)}\le 0\quad\forall\,w\in K_t.
}
Hence we must have $y\in\partial K_t$ and $\nu(x)$ is an outer normal of $K_t$ at $y$.

\emph{Step 5.} Let $L_t=K_t\cap\ov{\bbR^{n+1}_+}$.
We prove
\eq{
  \Sigma_t=\phi_t(\Sigma)\subset\partial L_t.
}

If  $x\in\op{int}(\Sigma)$, then by step 1, we have
\eq{
  y=\phi_t(x)\in\bbR^{n+1}_+.
}
Together with $y\in\partial K_t$ (by step 4) this implies $y\in\partial L_t\setminus P$.

If  $x\in\partial\Sigma$, then by step 1, $\phi_t(\partial\Sigma)\subset P$, so $y\in P$.
Let $x_j\in\op{int}(\Sigma)$ be any sequence with $x_j\to x$. Set $y_j:=\phi_t(x_j)$. By continuity, $y_j\to y$.
Since $y_j\in \partial L_t$ the limit point $y$ belongs to $\partial L_t\cap P$.

\emph{Step 6.}
We prove $ \ov{\partial L_t\setminus P}= \Sigma_t$. By steps 1, 2 and 5,
\eq{
  \op{int}(\Si_t)=\phi_t(\op{int}(\Si))\subset\partial L_t\setminus P\implies \Sigma_t \subset \ov{\partial L_t\setminus P}.
}
It remains to prove $\partial L_t\setminus P \subset \Sigma_t\setminus P$.


Let $y\in\partial L_t\setminus P$. Then $y\in\partial K_t$. Suppose $u\in\bbS^n$ is an outer unit normal to $K_t$ at $y$, i.e.
\eq{\label{eq:supporting-at-y}
  \ip{w-y}{u}\le 0\quad\forall\,w\in K_t.
}
We may write
\eq{\label{eq:y-decomposition}
  y=(x-t\cos\theta\,e_{n+1})+tb,
  \quad x\in \widehat{\Sigma},\quad b\in\bbB.
}
We claim that $b=u$, $x\in \partial \wh{\Si}$, and $u$ is an outer normal of $\widehat{\Sigma}$ at $x$. 

Indeed, take any $x_0\in \widehat{\Sigma}$ and any $c\in\bbB$, and set
\eq{
  w=(x_0-t\cos\theta\,e_{n+1})+tc\in K_t.
}
Plugging this $w$ and \eqref{eq:y-decomposition} into \eqref{eq:supporting-at-y} gives
\eq{
  0\ge \ip{w-y}{u}
  =\ip{x_0-x}{u}+t\ip{c-b}{u}.
}
With $x_0=x$ and $c=u$,
\eq{
  1\le \ip{b}{u}\implies b=u.
}
Now with $c=b=u$, the inequality becomes
\eq{
  0\ge \ip{x_0-x}{u}\quad\forall\,x_0\in \widehat{\Sigma},
}
so $x\in \partial \wh{\Si}$ and $u$ is an outer normal of $\widehat{\Sigma}$ at $x$.

Since $y_{n+1}>0$ and $t>0$, we have from \eqref{eq:y-decomposition} (with $b=u$)
\eq{
  y_{n+1}=x_{n+1}-t\cos\theta+tu_{n+1}>0.
}
This implies that $x_{n+1}>0$, $x\in \Si$ and $u=\nu(x)$ (otherwise, if $x_{n+1}=0$, then we would have $u_{n+1}\le \cos\theta$ and hence $y_{n+1}\le 0$).
Substituting $b=u=\nu(x)$ into \eqref{eq:y-decomposition} yields
\eq{
  y=x+t(\nu(x)-\cos\theta\,e_{n+1})=\phi_t(x)\in\Sigma_t\setminus P.
}

\emph{Step 7.} By the previous steps, $\Si_t$ is a strictly convex (i.e. the enclosed region $\wh{\Si_t}$ is a convex body and the second fundamental form of $\Si_t$ is positive definite), $\theta$-capillary hypersurface, and for each point $x\in \Sigma$, the outward unit normal at the point $\phi_t(x)\in\Sigma_t$ is $\nu(x)$. Let $\zeta=\nu(x)-\cos\theta\, e_{n+1}$. Then
\eq{
s_{\Sigma_t}(\zeta)&=\langle \phi_t(x),\nu(x)\rangle\\
&=\langle x+t(\nu(x)-\cos\theta e_{n+1}),\nu(x)\rangle\\
&=s_{\Sigma}(\zeta)+t\ell(\zeta),
}
Since $\Sigma_t$ has the same capillary support function as $\Sigma+t\Cth$, we conclude that
\eq{
  \Sigma_t=\Sigma+t\Cth.
}
\end{proof}

\end{document}
